\chapter{Introduction to Python for Developers}
\section{Introduction to Python}
\subsection{What is Python?}
\begin{itemize}
    \item General-purpose programming language
    \item Free and open source
\end{itemize}

\subsubsection{The print function}
\begin{itemize}
    \item Built-in Python function
    \item Displays code results
    \item Used for testing
    \item Can use either single or double quotations
\end{itemize}
\inputminted[
    firstline=1,
    lastline=4,
    highlightlines={1,3}
]{Python}{CodeFiles/IntroToPythonForDevelopers.py}

\subsubsection{Documenting with comments}
\begin{itemize}
    \item Comments explain code
    \item Starts with \hl{\#} symbol
\end{itemize}
\inputminted[
    firstline=6,
    lastline=8,
    highlightlines={6,7}
]{Python}{CodeFiles/IntroToPythonForDevelopers.py}

\subsection{Python variables, strings, and integers}
\subsubsection{Data types}
\begin{itemize}
    \item The kind of value
    \begin{itemize}
        \item Number
        \item Text
        \item True/False
    \end{itemize}
    \item Python assign data types automatically
\end{itemize}

\subsubsection{String data types and variables}
\begin{itemize}
    \item \textbf{Strings}: text values
    \item Can contain letters, numbers, and punctuation
    \item Use single or double quotes
    \item Double quotes for words with apostrophes
    \item use lowercase and underscores for spaces i.e. \hl{ingredient\_name}
\end{itemize}
\inputminted[
    firstline=10,
    lastline=10,
    highlightlines=10
]{Python}{CodeFiles/IntroToPythonForDevelopers.py}

\subsubsection{Integer data types and variables}
\begin{itemize}
    \item \textbf{Integer}: whole numbers
\end{itemize}
\inputminted[
    firstline=11,
    lastline=11,
    highlightlines=11
]{Python}{CodeFiles/IntroToPythonForDevelopers.py}

\subsubsection{Printing and adjusting variables}
\inputminted[
    firstline=13,
    lastline=19
]{Python}{CodeFiles/IntroToPythonForDevelopers.py}

\subsection{Common data types}
\subsubsection{Decimal data types}
\begin{itemize}
    \item \textbf{Floats}: numbers with decimal, e.g. 1.5
\end{itemize}

\section{Working with Data Types}
\subsection{Working with strings}
\subsection{Lists}
\subsection{Dictionaries}
\subsection{Sets and tuples}

\section{Control flows and loops}
\subsection{Conditional statements and operators}
\subsection{For loops}
\subsection{While loops}
\subsection{Building a workflow}